\section{Network Results}
\label{sec:network}

The classifier described in Section~\ref{sec:method} was deployed onto the measured
conductance ladder by post-training quantization, with no retraining. On the fifteen-level
ladder it reaches \num{0.8304} accuracy, \num{0.7536} macro-F1 and a Cohen's kappa of
\num{0.7417}, against \num{0.8571}, \num{0.7925} and \num{0.7819} in full precision.

That result is reproducible to four decimal places across a change that moved every
absolute current in the device export by a factor of \num{0.6338}, because each weight is
realized as $(G^+-G^-)/g_\mathrm{max}$ and a global scale cancels exactly. This was the
predicted outcome and it is the check that the area-factor correction of
Section~\ref{sec:norm} was a change of normalization and not of physics.

\subsection{Level placement dominates level count}
\label{sec:placement}

Sweeping the number of conductance levels from two to fifteen gives a curve that is
\emph{not} monotonic below eight levels. Two independent grid constructions were used ---
a subset of the stops the device actually reaches, and a geometrically spaced ideal grid
over the same range --- and they agree, so the non-monotonicity is not an artefact of how
the coarse grid was built.

\begin{table}[t]
\caption{Deployed accuracy against the number of conductance levels, for two grid
constructions.}
\label{tab:levels}
\centering
\begin{tabular}{@{}rcc@{}}
\toprule
levels & measured subset (acc / F1) & log-spaced (acc / F1) \\
\midrule
2  & 0.5625 / 0.2857 & 0.5625 / 0.2857 \\
3  & 0.5625 / 0.2857 & 0.5595 / 0.2823 \\
4  & 0.3631 / 0.2384 & 0.3423 / 0.2308 \\
5  & 0.5060 / 0.4714 & 0.4911 / 0.4541 \\
6  & 0.3512 / 0.1978 & 0.2292 / 0.2713 \\
8  & 0.8036 / 0.7155 & 0.7768 / 0.7002 \\
10 & 0.7619 / 0.6387 & 0.7976 / 0.7230 \\
12 & \textbf{0.8452 / 0.7628} & 0.7946 / 0.6995 \\
15 & \textbf{0.8304 / 0.7536} & 0.8274 / 0.7372 \\
\bottomrule
\end{tabular}
\end{table}

Redrawing the \emph{interior} levels at fixed count, keeping the endpoints, shows what is
happening. At eight levels, four different random placements gave accuracies from
\num{0.336} to \num{0.491}, while spacing the same eight evenly gave \num{0.804}. At six
levels the spread is \numrange{0.226}{0.777}. Level count alone does not determine
performance; \emph{where the levels sit} does.

The mechanism is the differential encoding on a ladder that spans four decades. Because
$|G_i-G_j|$ is dominated by $\max(i,j)$, the achievable weight set is approximately
$\{0,\pm G_k/g_\mathrm{max}\}$: dense near zero and very sparse near $\pm1$. A grid that
spends its levels at the bottom of the ladder buys almost no usable weight resolution, and
post-training quantization has no opportunity to compensate. With arbitrary placement the
classifier needs roughly ten to twelve levels to reach the software neighbourhood
reliably; the evenly spaced eight-level grid that scores \num{0.804} is a good placement,
not a typical one.

\subsection{Variability, and why write-verify is not optional}
\label{sec:variability}

Twenty perturbed devices were simulated at the corners of a box in interface trap density
($\pm\SI{20}{\percent}$), fixed interface charge ($\pm\SI{10}{\percent}$) and ferroelectric
thickness ($\pm\SI{0.3}{\nano\metre}$). These are corners --- at least two axes off-nominal
at once --- so they are a worst-case envelope and not a Gaussian distribution, and they are
never injected as one. They are used the only way they legitimately can be: the locked
network is deployed onto each corner device under three levels of per-device calibration.

\begin{table}[t]
\caption{Deployed accuracy across the twenty corner devices, by how much per-device
calibration is allowed.}
\label{tab:calib}
\centering
\begin{tabular}{@{}lcc@{}}
\toprule
calibration & median accuracy & range \\
\midrule
none & \num{0.2426} & \numrange{0.069}{0.771} \\
one gain scalar per device & \num{0.7768} & \numrange{0.244}{0.830} \\
full write-verify per level & \num{0.8304} & \numrange{0.798}{0.860} \\
\bottomrule
\end{tabular}
\end{table}

The full-calibration median equals the nominal device's accuracy to four decimal places,
and every corner lands between \num{0.798} and \num{0.860}. A single gain trim --- a
transimpedance adjustment, essentially free --- recovers about two thirds of the gap; the
remaining \num{0.054} needs per-level placement. This is why Section~\ref{sec:levels}
reports the fifteen-level claim as a write-verify claim without treating it as a
concession: the array needs write-verify to work at all.

\subsection{Overlap does not predict accuracy; misplacement does}

Decomposing each corner device's ladder into a rigid shift in $\log G$ plus a residual
shape error gives median shifts of \num{0.51} decades and a median residual scatter of
\num{0.178} decades. The shift is what a gain trim removes. What remains is level
misplacement, and that is what accuracy tracks: against the worst single-level residual,
gain-trimmed accuracy has a Spearman correlation of \num{-0.85}, compared with
\num{-0.79} against the residual standard deviation and \num{-0.71} against the mean
shift. Accuracy holds between \num{0.78} and \num{0.83} out to roughly \num{0.6} decades
of worst-level misplacement and collapses to about \num{0.25} by \num{0.9}. Full
write-verify is flat against the same axis, as it must be.

This produces a design rule that is more useful than a level count: \textbf{if only a gain
trim is affordable, interface-charge control must hold worst-level placement inside about
\num{0.6} decades.} Combined with the axis decomposition below, that is a specification a
process team can act on.

The corollary matters for how the ensemble is reported. After a per-device gain trim,
thirteen of fourteen adjacent level pairs still overlap across the ensemble --- half-spread
\num{0.36} decades against a spacing of \num{0.28} --- and yet several of those corners
deliver nominal accuracy. Ensemble overlap says that a \emph{shared} grid cannot uniquely
address a level. It does not say the deployed weight is wrong enough to matter.

\subsection{Interface charge, not ferroelectric thickness}
\label{sec:interface}

Decomposing the corner spread by axis, at the eighth ladder stop:

\begin{table}[t]
\caption{Spread of group medians at LTP level 8, by perturbed parameter.}
\label{tab:axes}
\centering
\begin{tabular}{@{}llc@{}}
\toprule
parameter & perturbation & spread of group medians \\
\midrule
fixed interface charge & $\pm\SI{10}{\percent}$ & \textbf{\num{1.31} decades} \\
interface trap density & $\pm\SI{20}{\percent}$ & \num{0.85} decades \\
ferroelectric thickness & $\pm\SI{0.3}{\nano\metre}$ & \textbf{\num{0.07} decades} \\
\bottomrule
\end{tabular}
\end{table}

This is the opposite of the expected ordering. Ferroelectric thickness sets the coercive
voltage, so it looks like the obvious lever, and it turns out to be almost irrelevant to
where the analog levels land: a \SI{0.3}{\nano\metre} thickness tolerance costs
\num{0.07} decades, while \SI{10}{\percent} of fixed interface charge costs \num{1.31}. An
order-of-magnitude check agrees --- a charge change of
\SI{0.7e12}{\per\square\centi\metre} across a \SI{1}{\nano\metre} oxide is about
\SI{32}{\milli\volt} of threshold shift, roughly half a decade at
\SI{63}{\milli\volt\per\text{dec}}, rising once the ferroelectric divider is included.

\textbf{Analog level placement in this device is an interface-charge control problem, not a
ferroelectric thickness control problem.} That is the most directly actionable result here.

\subsection{Cycle-to-cycle noise costs essentially nothing}

Injecting the measured cycle-to-cycle noise per conductance branch, the fifteen-level
deployment loses \num{0.001} accuracy: \num{0.8304} becomes \num{0.8294}. The knee is at
about \numrange{0.1}{0.2} decades of injected noise, so the device sits five to nine times
below it, and the conclusion survives the \SI{22}{\percent} uncertainty in the measured
standard deviation --- even at \num{0.05} decades, roughly two standard errors above the
measurement, accuracy is still \num{0.8175}.

Quantizing instead to the eight levels that pass the three-sigma separability test costs
\num{0.036} accuracy. That is the wrong trade, and the reason is worth stating because it
is a design-methodology point rather than a device one. The separability criterion returns
an \emph{integer}, derived from a three-sigma test on a standard deviation that is itself
known to about \SI{22}{\percent}. An earlier pass on a partial download of the same run
saw four complete repeats, put the standard deviation at \num{0.0358} and the usable count
at six, and the accuracy penalty at \numrange{0.42}{0.48}. Six levels lands in the
catastrophic regime of Table~\ref{tab:levels}; eight lands on its good-placement point. The
cost of following the criterion therefore swung by an order of magnitude on which integer
the test happened to return. A design rule that unstable is not one to quantize a network
by. The robust statement is the direct one: at the measured noise, keeping all fifteen
levels costs \num{0.001}, so keep them.

\subsection{Distinguishability is the wrong figure of merit}

Three independent analyses in this work measure whether conductance levels can be told
apart on readout --- the corner ensemble's level overlap, the overlap remaining after a
per-device gain trim, and the three-sigma separability of the cycle-to-cycle run --- and
\textbf{none of them predicts accuracy.} What predicts accuracy is how far the deployed
weight ends up from the weight that was wanted: per-device level misplacement, and where
the levels sit.

The reason is that separability asks whether a level can be distinguished from its
neighbour on \emph{readout}, while the network only requires the realized weight to sit
near its target. Levels that overlap under noise remain monotonic and still carry weight
information. Overlap and separability describe readout addressing; only deployed-weight
error moves the classifier.

One detail should not be over-read: deep in the collapse, at \num{0.2} decades of injected
noise, the eight-level grid slightly beats the fifteen-level one (\num{0.4415} against
\num{0.3919}). A coarser grid is more noise-robust once noise exceeds level spacing, but
both are far below usable accuracy there, so it is not a design argument.

\subsection{Robustness}

Deployed accuracy under the measured device conditions:

\begin{table}[t]
\caption{Deployed accuracy under measured and stress conditions.}
\label{tab:robust}
\centering
\begin{tabular}{@{}lccc@{}}
\toprule
condition & acc & macro-F1 & $\kappa$ \\
\midrule
full precision & 0.857 & 0.793 & 0.782 \\
nominal, \SI{300}{\kelvin} & 0.830 & 0.754 & 0.742 \\
\SI{250}{\kelvin}, measured ladder & 0.821 & 0.739 & 0.729 \\
\SI{350}{\kelvin}, measured ladder & 0.839 & 0.766 & 0.753 \\
retention $-\SI{0.96}{\percent}$, measured & 0.821 & 0.721 & 0.728 \\
retention $-\SI{10}{\percent}$, stress & 0.819 & 0.740 & 0.725 \\
\bottomrule
\end{tabular}
\end{table}

Two things about Table~\ref{tab:robust} deserve saying rather than leaving for a reader to
find. The measured retention loss scores marginally \emph{below} the ten-percent stress
case. Both sit about \num{0.03} macro-F1 under nominal and the ordering is not meaningful:
a retention factor is a uniform multiplicative gain applied after quantization, and a
spiking classifier with fixed thresholds responds to a small uniform gain in discrete
spike-count steps rather than smoothly. Individual conditions also moved by at most four
beats out of \num{336} between the pre- and post-correction runs, which is float32 rounding
at quantization boundaries --- the levels are identical after normalization but not
bit-identical, so a handful of weights land on the other side of a bucket edge. Worth
knowing that the deployment is boundary-sensitive at the one-beat level; not worth reading
anything into.

\subsection{Energy, and the scope of the claim}
\label{sec:energy}

Every count in the energy budget was taken from the deployed model rather than assumed.
The network has \num{261440} weights and therefore \num{522880} devices in differential
pairs. Measured activity is \num{196.1} input spikes per beat
(\SI{5.86}{\percent} line activity) and \num{3455.7} hidden spikes per beat
(\SI{1.94}{\percent}, \num{21.6} spikes per neuron).

\begin{table}[t]
\caption{Energy per heartbeat, all counts measured on the deployed network.}
\label{tab:energy}
\centering
\begin{tabular}{@{}lr@{}}
\toprule
component & energy \\
\midrule
device write, whole array, one time & \SI{69.6}{\pico\joule} \\
device read, event driven & \SI{103.9}{\pico\joule} / beat \\
device read, fully clocked & \SI{4.62}{\nano\joule} / beat \\
neuron membrane & \SI{432.0}{\pico\joule} / beat \\
\textbf{core compute} & \textbf{\SI{535.8}{\pico\joule} / beat} \\
peripheral ADC (assumed \SI{1}{\pico\joule}/conversion) & \SI{183.0}{\nano\joule} / beat \\
\bottomrule
\end{tabular}
\end{table}

Two results follow. First, \textbf{programming the entire array costs \num{0.13} of a
single inference}, and non-volatility means it is paid once and never again. That is the
concrete argument for a ferroelectric weight, and it is a measured number rather than an
appeal to the physics.

Second, and more important, \textbf{the assumed analog-to-digital conversion term is
\num{342} times the core compute.} On this workload the ferroelectric array is not the
limiting element of a complete system. Reporting only the core figure would overclaim by
more than two orders of magnitude, so the peripheral cost is a bar in the figure rather
than a sentence in a caveat. The \SI{1}{\pico\joule} per conversion figure is the single
assumption in this budget and it is labelled as such on the axis; the conclusion is not
sensitive to it, since the term would have to fall by two orders of magnitude before the
array became the bottleneck.
